% Parte B: erros, observações, exemplos, questões, respostas
\ErrosComuns

\begin{erro}{Achar que antecedente falso torna o condicional falso}{3-E1}
\campo{O erro.} Classificar \enquote{se $3\cdot4=10$, então $1>2$} como falso \enquote{porque $3\cdot4$ não é 10}.
\campo{Por que parece plausível.} Lê-se o condicional como se afirmasse o antecedente. No dia a dia, quem diz \enquote{se\ldots} costuma achar que a condição pode acontecer.
\campo{Onde quebra.} No teorema \enquote{para todo inteiro $n$, se $n$ é múltiplo de 4, então $n$ é par}. Para $n=6$, o antecedente é falso. Se isso tornasse o condicional falso, o teorema seria falso, e quase nenhum teorema sobreviveria.
\campo{Como detectar.} Antes de classificar, olhe o antecedente. Se ele é F, a resposta já está dada, e você não precisa nem ler o consequente.
\campo{Como corrigir.} Lembre a promessa (construção 3-H1): se a condição não acontece, a promessa não é quebrada. Antecedente F: condicional V.
\campo{Onde aparece.} Exemplo \ref{ex:3.1}; questões \ref{u3:q1}, \ref{u3:q8} e \ref{u3:q13}.
\end{erro}

\begin{erro}{Confundir recíproca ou inversa com contrapositiva}{3-E2}
\campo{O erro.} De \enquote{se é quadrado, tem quatro ângulos retos}, concluir \enquote{se tem quatro ângulos retos, é quadrado} (recíproca) ou \enquote{se não é quadrado, não tem quatro ângulos retos} (inversa).
\campo{Por que parece plausível.} As três frases têm as mesmas palavras; trocar a ordem ou acrescentar \enquote{não} parece inofensivo. Muitas vezes a recíproca também é verdadeira, e o estudante generaliza.
\campo{Onde quebra.} Num retângulo que não é quadrado: tem quatro ângulos retos e não é quadrado. A recíproca e a inversa falham; a frase original, não.
\campo{Como detectar.} Compare as posições: você trocou a ordem \emph{e} negou as duas partes? Só isso é a contrapositiva. Se fez só uma das duas coisas, procure um contraexemplo.
\campo{Como corrigir.} Use a construção 3-H3: $p\rightarrow q$ equivale a ${\nao}q\rightarrow{\nao}p$, e só a ela.
\campo{Onde aparece.} Exemplo \ref{ex:3.6}; questões \ref{u3:q4} e \ref{u3:q10}.
\end{erro}

\begin{erro}{Inverter \enquote{necessária}, \enquote{suficiente} e \enquote{somente se}}{3-E3}
\campo{O erro.} Ler \enquote{ser par é condição necessária para ser múltiplo de 6} como \enquote{se é par, então é múltiplo de 6}; ler \enquote{$p$ somente se $q$} como $q\rightarrow p$.
\campo{Por que parece plausível.} \enquote{Necessário} soa como \enquote{o que vem primeiro}, e \enquote{somente se} parece reforçar o \enquote{se}.
\campo{Onde quebra.} Em $n=2$: é par e não é múltiplo de 6. A leitura errada produz uma frase falsa a partir de uma verdadeira.
\campo{Como detectar.} Teste a sua tradução com um número em que a condição necessária vale e a outra não. Se a tradução fica falsa e o enunciado era verdadeiro, ela está invertida.
\campo{Como corrigir.} A condição \emph{suficiente} fica na frente da seta; a \emph{necessária}, depois. \enquote{$p$ somente se $q$} é $p\rightarrow q$.
\campo{Onde aparece.} Exemplo \ref{ex:3.2}; questão \ref{u3:q3}.
\end{erro}

\begin{erro}{Ler o \enquote{se, e somente se} como \enquote{e}}{3-E4}
\campo{O erro.} Classificar como falso um bicondicional em que as duas partes são falsas, como \enquote{$5>9\leftrightarrow 2=3$}.
\campo{Por que parece plausível.} Pensa-se que \enquote{$p$ se, e somente se, $q$} afirma $p$ e afirma $q$.
\campo{Onde quebra.} Na definição \enquote{$n$ é primo se, e somente se, tem exatamente dois divisores naturais}. Para $n=1$, as duas partes são falsas; se o bicondicional fosse falso aí, a definição de primo estaria errada.
\campo{Como detectar.} Olhe a linha F e F. Se você marcou F nela, leu um \enquote{e}.
\campo{Como corrigir.} O bicondicional não afirma as partes: afirma que elas \emph{andam juntas}. Mesmo valor, V; valores diferentes, F.
\campo{Onde aparece.} Exemplo \ref{ex:3.3}; questão \ref{u3:q5}.
\end{erro}

\begin{erro}{Exigir causalidade}{3-E5}
\campo{O erro.} Dizer que \enquote{se 2 é par, então Natal fica no RN} \enquote{não faz sentido} ou é falso, porque uma coisa não causa a outra.
\campo{Por que parece plausível.} Na língua, \enquote{se\ldots\ então} costuma indicar causa ou razão; os critérios antigos que pediam uma ligação entre as partes, como o da \enquote{coesão} (observação 3-O1), ainda soam naturais.
\campo{Onde quebra.} Nos próprios teoremas: \enquote{ser múltiplo de 4} não \emph{causa} \enquote{ser par}, e \enquote{se $n$ é múltiplo de 4, então $n$ é par} é verdadeiro. A matemática não trata de causas, e sim de valores.
\campo{Como detectar.} Se você está perguntando \enquote{por que uma coisa leva à outra?} em vez de \enquote{quais são os valores?}, saiu do conectivo $\rightarrow$.
\campo{Como corrigir.} O conectivo $\rightarrow$ segue Fílon: só importam os valores. V e V dão V, haja ou não relação entre as partes.
\campo{Onde aparece.} Exemplo \ref{ex:3.1}; questão \ref{u3:q1}.
\end{erro}

% ------------------------------------------------------------
\Observacoes

\begin{obs}{3-O1}
\textbf{Por que o critério de Fílon venceu.} O critério de Fílon tem consequências estranhas: \enquote{se Natal fica em São Paulo, então a Lua é de queijo} é verdadeiro, só porque o antecedente é falso. Para Diodoro, seria falso: é \emph{possível} que o antecedente fosse verdadeiro e o consequente falso. Foi esse tipo de caso, o de antecedente falso por acaso, que incomodou Diodoro e, vinte e dois séculos depois, C.~I.~Lewis.\src{sepDialectical}\src{organonF} Por que, então, a matemática o adotou?

Sexto Empírico registra, além dos de Fílon e Diodoro, dois critérios mais exigentes. Pelo da \enquote{coesão}, o condicional é correto quando a negação do consequente é incompatível com o antecedente. Pelo da \enquote{inclusão}, o consequente precisa estar contido no antecedente, e aí até \enquote{se é dia, é dia} seria falso.\src{ls35b}

A primeira razão para a escolha da matemática é que o critério de Fílon é o único que depende \emph{apenas} dos valores das partes. Os de Diodoro, da coesão e da inclusão perguntam pelo que é \emph{possível} ou pelo que está \enquote{contido} numa frase — perguntas que, lidas hoje, uma tabela não responde. A segunda razão aparece nos teoremas. Quando dizemos \enquote{para todo inteiro $n$, se $n$ é múltiplo de 4, então $n$ é par}, queremos que a frase seja verdadeira para $n=6$ também, em que o antecedente é falso. Se o antecedente falso tornasse o condicional falso — ou \enquote{sem valor} —, quase nenhum teorema seria verdadeiro. A regra de Fílon é exatamente a que faz os teoremas funcionarem: um \enquote{se} geral só é desmentido por um caso em que a hipótese vale e a tese não, ou seja, por um contraexemplo.

O preço são as frases estranhas com antecedente falso. A lógica as aceita como verdadeiras e \emph{inofensivas}: elas nunca permitem concluir nada falso a partir de premissas verdadeiras. Quem quiser um \enquote{se} mais forte — com necessidade, causa ou relevância — usa outras lógicas, como a implicação estrita de Lewis.\src{openlogicPar}
\end{obs}

\begin{obs}{3-O2}
\textbf{As cartas de Wason com conteúdo social.} Mude apenas a história. Você trabalha num bar e deve fazer cumprir a regra \enquote{\emph{se uma pessoa bebe cerveja, então tem pelo menos 18 anos}}. Quatro cartões mostram, de um lado, a bebida e, do outro, a idade de quatro clientes: \enquote{cerveja}, \enquote{refrigerante}, \enquote{25 anos}, \enquote{16 anos}. Quais você vira? A maioria responde certo, e depressa: \enquote{cerveja} e \enquote{16 anos}. Em 1982 e 1983, os psicólogos Richard Griggs e James Cox mostraram que, com essa versão, a maioria dos participantes acerta.\src{griggsCox1982}\src{sussexWason}

A estrutura lógica é idêntica à das letras e números: $p$ = \enquote{bebe cerveja}, $q$ = \enquote{tem 18 anos ou mais}, e as cartas a virar são $p$ e ${\nao}q$. Por que, então, a diferença? Em 1989, Leda Cosmides propôs que os seres humanos raciocinam especialmente bem quando a regra é um \enquote{contrato social} e o erro seria deixar passar um trapaceiro.\src{cosmides1989} A explicação é disputada: outros pesquisadores argumentam que o efeito vem de como o problema é entendido, e não de uma habilidade especial.\src{sperber} Para nós, fica a lição prática: diante de um condicional abstrato, traduza-o para uma regra que alguém poderia \emph{violar} e procure o violador.
\end{obs}

\begin{obs}{3-O3}
\textbf{Corvos na lógica.} A disputa entre Fílon e Diodoro ficou famosa a ponto de virar piada: o poeta alexandrino Calímaco escreveu um epigrama em que os corvos, pousados nos telhados, perguntam uns aos outros \enquote{o que se segue de quê?}.\src{whiteCrows} Outros corvos voltaram no século XX. Em 1945, o filósofo Carl Hempel examinou a frase \enquote{todos os corvos são pretos}, que é um condicional: \enquote{se algo é corvo, então é preto}. Pela contrapositiva, ela equivale a \enquote{se algo não é preto, então não é corvo}.\src{hempel1945}\src{ravenWiki} Observar um corvo preto parece apoiar a primeira forma. Mas, então, observar uma maçã verde — algo que não é preto e não é corvo — apoiaria a segunda, e portanto a primeira. Uma maçã verde seria evidência de que os corvos são pretos?

O \enquote{paradoxo dos corvos} não está na lógica: as duas formas são de fato equivalentes, como mostra a construção 3-H3. O problema está em como observações \emph{confirmam} uma lei geral, e gerou uma longa discussão sobre indução e probabilidade.\src{fitelson} Há uma leitura intuitiva: como existem muitíssimo mais coisas não pretas do que corvos, uma maçã verde confirma a lei, mas de forma quase nula. Para a matemática, que prova em vez de colecionar exemplos, a lição é outra: para refutar \enquote{todo corvo é preto}, basta um corvo branco — o caso \enquote{antecedente V, consequente F}.
\end{obs}

\begin{obs}{3-O4}
\textbf{O \enquote{se} nos computadores.} O \texttt{if} das linguagens de programação não é o conectivo $\rightarrow$: é uma instrução de controle (\enquote{se a condição for verdadeira, execute isto}), que não tem valor verdadeiro ou falso. Mas a ideia de Fílon aparece nos testes de software: a regra \enquote{se o pedido é internacional, então a taxa é cobrada} só é violada por um pedido internacional sem taxa, e os testadores procuram exatamente esse caso, como nas cartas de Wason. Já o bicondicional entre dois valores booleanos é a igualdade: em Python, \texttt{p == q}; nos circuitos, a porta XNOR, o contrário da XOR da Unidade 2.
\end{obs}

% ------------------------------------------------------------
\Exemplos

\begin{exemplo}{3.1}{1}{adaptado do FME, item 5}{3-H1, 3-E1, 3-E5}\label{ex:3.1}
Classifique em verdadeira ou falsa.
\begin{multicols}{2}
\begin{enumerate}[label=\alph*)]
  \item Se $3\mid 12$, então $3\mid 24$.
  \item Se $2+2=4$, então 2 é ímpar.
  \item Se $1>2$, então 2 é primo.
  \item Se 7 é par, então $7=0$.
  \item Se Natal fica no RN, então 9 é um quadrado perfeito.
  \item Se $\sqrt2$ é racional, então $\pi=3$.
\end{enumerate}
\end{multicols}
\plano Calcular o valor do antecedente e do consequente; aplicar a tabela: o único caso F é V seguido de F.
\resolucao
\begin{center}\small
\begin{tabular}{lccc}\toprule
& antecedente & consequente & $p\rightarrow q$\\\midrule
a) & V & V ($24=3\cdot8$) & \textbf{V}\\\fino
b) & V & F & \textbf{F}\\\fino
c) & F & V & \textbf{V}\\\fino
d) & F & F & \textbf{V}\\\fino
e) & V & V ($9=3^2$) & \textbf{V}\\\fino
f) & F & F & \textbf{V}\\\bottomrule
\end{tabular}
\end{center}
\comentario Só (b) é falso. Em (c), (d) e (f), o antecedente é falso e o condicional é verdadeiro (erro 3-E1); (d) e (f) parecem absurdos, mas são inofensivos (veja a observação 3-O1). Em (e), nada liga Natal ao número 9: o condicional é verdadeiro porque as duas partes são verdadeiras (erro 3-E5).
\end{exemplo}

\begin{exemplo}{3.2}{1}{autoral}{3-U1, 3-E3}\label{ex:3.2}
Considere, para inteiros $n$, a proposição \enquote{se $n$ é múltiplo de 4, então $n$ é par}.
\begin{enumerate}[label=\alph*)]
  \item Escreva-a com as palavras \enquote{suficiente}, \enquote{necessária} e \enquote{somente se}.
  \item Decida se é verdadeira: \enquote{ser par é condição suficiente para ser múltiplo de 4}.
  \item Para números reais $x$, \enquote{$x^2=9$} é condição necessária ou suficiente para \enquote{$x=3$}?
\end{enumerate}
\plano Antecedente = suficiente; consequente = necessária; \enquote{$p$ somente se $q$} = $p\rightarrow q$. Testar com números.
\resolucao
a) \enquote{Ser múltiplo de 4 é condição \emph{suficiente} para ser par}; \enquote{ser par é condição \emph{necessária} para ser múltiplo de 4}; \enquote{$n$ é múltiplo de 4 \emph{somente se} $n$ é par}.
\par b) A frase diz \enquote{se $n$ é par, então $n$ é múltiplo de 4} — a recíproca. \textbf{Falsa}: $n=6$ é par e não é múltiplo de 4.
\par c) \enquote{Se $x=3$, então $x^2=9$} é verdadeiro; logo, $x^2=9$ é \textbf{necessária} para $x=3$. Não é suficiente: $x=-3$ satisfaz $x^2=9$ sem ser 3.
\comentario Uma condição pode ser necessária sem ser suficiente, e vice-versa. Quando é as duas coisas, temos um bicondicional — o tema do próximo exemplo.
\end{exemplo}

\begin{exemplo}{3.3}{2}{autoral}{3-H2, 3-U5, 3-E4}\label{ex:3.3}
\begin{enumerate}[label=\alph*)]
  \item Classifique: (i) $2\cdot6=12\leftrightarrow 2\cdot6+7=19$;\ (ii) $5>9\leftrightarrow 2=3$;\ (iii) $2\mid7\leftrightarrow 7$ é primo;\ (iv) $3^2=9\leftrightarrow 3$ é par.
  \item Para todo inteiro $n$, vale \enquote{$n$ é par $\leftrightarrow$ $n^2$ é par}? E \enquote{$n$ é múltiplo de 6 $\leftrightarrow$ $n$ é múltiplo de 3}?
  \item Por que a definição \enquote{um número é par se $n=2k$, com $k$ inteiro} deve ser lida como \enquote{se, e somente se}?
\end{enumerate}
\plano Em (a), comparar os valores. Em (b), um bicondicional geral exige os dois sentidos; um contraexemplo em qualquer sentido o derruba.
\resolucao
a) (i) V e V: \textbf{V}. (ii) F e F: \textbf{V}. (iii) F e V: \textbf{F}. (iv) V e F: \textbf{F}.
\par b) Primeiro: se $n=2k$, então $n^2=2(2k^2)$ é par; e, se $n$ é ímpar, $n^2$ é ímpar (veja \enquote{Onde é usado}), logo $n^2$ par obriga $n$ par. \textbf{Vale nos dois sentidos.} Segundo: múltiplo de 6 $\rightarrow$ múltiplo de 3 vale, mas a volta falha em $n=3$. \textbf{Não vale.}
\par c) Se a definição dissesse só \enquote{se}, ela não excluiria que outros números também fossem chamados de pares. Uma definição precisa dizer exatamente quais objetos recebem o nome: é um bicondicional.
\comentario O item (ii) é a armadilha do erro 3-E4: duas partes falsas dão bicondicional \emph{verdadeiro}. O \enquote{se, e somente se} afirma que as partes andam juntas, não que valem.
\end{exemplo}

\begin{exemplo}{3.4}{2}{autoral}{3-U3}\label{ex:3.4}
Sabe-se que a proposição $(p\wedge q)\rightarrow(r\vee{\nao}s)$ é \textbf{falsa}. Determine os valores de $p$, $q$, $r$, $s$ e, em seguida, classifique:\quad a) $r\rightarrow p$\quad b) $p\leftrightarrow s$\quad c) $q\rightarrow r$\quad d) ${\nao}p\vee s$.
\plano Um condicional falso tem antecedente V e consequente F. Desmontar cada parte com as tabelas de $\wedge$ e $\vee$.
\resolucao
O antecedente $p\wedge q$ é V; logo, $p$ e $q$ são V. O consequente $r\vee{\nao}s$ é F; logo, $r$ é F e ${\nao}s$ é F, isto é, $s$ é V.
\par a) F $\rightarrow$ V: \textbf{V}. b) V $\leftrightarrow$ V: \textbf{V}. c) V $\rightarrow$ F: \textbf{F}. d) F $\vee$ V: \textbf{V}.
\comentario Saber que um condicional é \emph{verdadeiro} quase não informa (há três casos); saber que é \emph{falso} fixa tudo (há um caso só). É por isso que os concursos gostam do formato \enquote{sabe-se que a afirmação\ldots\ é falsa}.
\end{exemplo}

\begin{exemplo}{3.5}{2}{OBMEP 2007, 1ª fase, nível 2, q.~8}{3-U2, 3-U4, 3-H3}\label{ex:3.5}
A mãe de César deu a ele as seguintes instruções para fazer um bolo:
\begin{itemize}
  \item se colocar ovos, não coloque creme;
  \item se colocar leite, não coloque laranja;
  \item se não colocar creme, não coloque leite.
\end{itemize}
Seguindo essas instruções, César pode fazer um bolo com:
\begin{enumerate}[label=(\alph*)]
  \item ovos e leite, mas sem creme.
  \item creme, laranja e leite, mas sem ovos.
  \item ovos e creme, mas sem laranja.
  \item ovos e laranja, mas sem leite e sem creme.
  \item leite e laranja, mas sem creme.
\end{enumerate}
\plano Reescrever a terceira regra pela contrapositiva e testar cada alternativa: uma regra é violada quando o antecedente acontece e o consequente não.
\resolucao
A terceira regra, ${\nao}C\rightarrow{\nao}L$, equivale a $L\rightarrow C$: \enquote{se colocar leite, coloque creme}.
\begin{itemize}
  \item (a) tem leite sem creme: viola a 3ª regra.
  \item (b) tem leite e laranja: viola a 2ª.
  \item (c) tem ovos e creme: viola a 1ª.
  \item (e) tem leite e laranja: viola a 2ª (e a 3ª).
  \item (d) ovos sem creme: 1ª ok; sem leite: a 2ª tem antecedente falso, e a 3ª (sem creme $\rightarrow$ sem leite) tem consequente verdadeiro, ok. \textbf{Alternativa (d).}
\end{itemize}
\comentario Em (d), a 2ª regra e a 3ª, escrita como \enquote{se leite, então creme}, são satisfeitas \enquote{por vacuidade}: como não há leite, nada exigem. Na forma em que a mãe a disse (\enquote{se não puser creme, não ponha leite}), a 3ª tem antecedente V e consequente V. Ser verdadeiro \enquote{por vacuidade} depende da forma em que o condicional está escrito. É o caso \enquote{antecedente F} da tabela — e mostra que o critério de Fílon é o que torna as instruções da mãe sensatas.
\end{exemplo}

\begin{exemplo}{3.6}{3}{histórico (Wason, 1966)}{3-E2, 3-U4, 3-O3}\label{ex:3.6}
Volte às quatro cartas do prefácio: \fbox{A}\ \fbox{K}\ \fbox{4}\ \fbox{7}, com a regra \enquote{se uma carta tem vogal de um lado, então tem número par do outro}. Quais cartas é preciso virar, e só essas, para saber se a regra vale para as quatro? Justifique a escolha de cada carta.
\plano A contrapositiva aparece aqui num uso novo: decidir o que verificar. Uma carta só desmente a regra se tiver vogal \emph{e} número ímpar (o caso F do condicional). Para cada carta visível, perguntar: o verso pode produzir esse caso?
\resolucao
\begin{itemize}
  \item \textbf{A} (vogal, antecedente V): se o verso for ímpar, a regra cai. \textbf{Virar.}
  \item \textbf{K} (consoante, antecedente F): a regra não diz nada sobre consoantes. Não virar.
  \item \textbf{4} (par, consequente V): qualquer letra no verso deixa a regra intacta — se for vogal, ótimo; se for consoante, a regra não fala dela. Não virar.
  \item \textbf{7} (ímpar, consequente F): se o verso for vogal, a regra cai. \textbf{Virar.} É a contrapositiva: \enquote{se é ímpar, então não é vogal}.
\end{itemize}
Resposta: \textbf{A e 7}.
\comentario A escolha errada mais comum, \enquote{A e 4}, procura confirmações: vira as cartas que \enquote{combinam} com a regra e esquece a contrapositiva (erro 3-E2). Em 1971, só 2 de 16 estudantes acertaram a versão abstrata;\src{wasonShapiro} com conteúdo social, a maioria acerta (3-O2). E o 7 é a \enquote{coisa não preta} de Hempel (3-O3): virá-lo é verificar se ela esconde um corvo, isto é, uma vogal.
\end{exemplo}

% ------------------------------------------------------------
\Questoes

\begin{questao}{1}{FME, exercício 4}{3-H1, 3-H2, 3-E1, 3-E5}\label{u3:q1}
Classifique em verdadeira ou falsa cada uma das proposições abaixo.
\begin{multicols}{2}
\begin{enumerate}[label=\alph*)]
  \item $2-1=1\rightarrow 5+7=3\cdot4$
  \item $2^2=4\leftrightarrow(-2)^2=4$
  \item $5+7\cdot1=10\rightarrow 3\cdot3=9$
  \item $\operatorname{mdc}(3,6)=1\leftrightarrow 4$ é número primo
  \item $2\mid 8\rightarrow\operatorname{mmc}(2,8)=2$
  \item $6\le2\leftrightarrow 6-2\ge0$
  \item $\tfrac35<\tfrac27\rightarrow 3\cdot7=2\cdot5$
\end{enumerate}
\end{multicols}
\resp{a) V\ \ b) V\ \ c) V\ \ d) V\ \ e) F\ \ f) F\ \ g) V}
\end{questao}

\begin{questao}{1}{FME, exercício 5}{3-H1, 3-H2}\label{u3:q2}
Admitindo que $p$ e $q$ são verdadeiras e $r$ é falsa, determine o valor (V ou F) de cada proposição abaixo.
\begin{multicols}{2}
\begin{enumerate}[label=\alph*)]
  \item $p\rightarrow r$
  \item $p\leftrightarrow q$
  \item $r\rightarrow p$
  \item $(p\vee r)\leftrightarrow q$
  \item $p\rightarrow(q\rightarrow r)$
  \item $p\rightarrow(q\vee r)$
  \item ${\nao}p\leftrightarrow{\nao}q$
  \item ${\nao}p\leftrightarrow r$
\end{enumerate}
\end{multicols}
\resp{a) F\ \ b) V\ \ c) V\ \ d) V\ \ e) F\ \ f) V\ \ g) V\ \ h) V}
\end{questao}

\begin{questao}{1}{autoral}{3-U1, 3-E3}\label{u3:q3}
Para inteiros $n$, considere a proposição verdadeira \enquote{se $n$ é múltiplo de 6, então $n$ é par}. Qual alternativa diz a mesma coisa?
\begin{enumerate}[label=(\alph*)]
  \item Ser par é condição suficiente para ser múltiplo de 6.
  \item Ser múltiplo de 6 é condição necessária para ser par.
  \item Ser par é condição necessária para ser múltiplo de 6.
  \item $n$ é par somente se $n$ é múltiplo de 6.
  \item $n$ é múltiplo de 6 se $n$ é par.
\end{enumerate}
\resp{(c)}
\end{questao}

\begin{questao}{1}{autoral}{3-U4, 3-E2}\label{u3:q4}
Considere: \enquote{se um quadrilátero é um quadrado, então ele tem quatro ângulos retos}. Escreva a recíproca, a inversa e a contrapositiva, e diga quais delas são verdadeiras para todo quadrilátero.
\resp{Recíproca: se tem quatro ângulos retos, é quadrado (F). Inversa: se não é quadrado, não tem quatro ângulos retos (F). Contrapositiva: se não tem quatro ângulos retos, não é quadrado (V).}
\end{questao}

\begin{questao}{2}{autoral}{3-U5, 3-E4}\label{u3:q5}
Para cada $x\in\{-2,-1,0,1,2\}$, decida o valor lógico de \enquote{$x^2=1\leftrightarrow x=1$} e de \enquote{$x^2=1\wedge x=1$}. Para quantos valores de $x$ cada uma é verdadeira? Explique a diferença. Depois, escreva a definição \enquote{um número natural é primo quando tem exatamente dois divisores naturais} como um bicondicional e dê o valor de cada parte, e do bicondicional, para o número 1.
\resp{$\leftrightarrow$: V para $x=-2,0,1,2$ (4 valores);\ \ $\wedge$: V só para $x=1$ (1 valor);\ \ \enquote{$n$ é primo $\leftrightarrow$ $n$ tem exatamente dois divisores naturais}: para 1, as duas partes são F e o bicondicional é V.}
\end{questao}

\begin{questao}{2}{FME, exercício 6}{3-U3}\label{u3:q6}
Sendo a proposição $p\rightarrow(r\vee s)$ falsa e a proposição $(q\wedge{\nao}s)\leftrightarrow p$ verdadeira, classifique em verdadeira ou falsa as afirmações $p$, $q$, $r$ e $s$.
\resp{$p$: V;\ \ $q$: V;\ \ $r$: F;\ \ $s$: F}
\end{questao}

\begin{questao}[integradora]{2}{autoral}{3-U1, 3-U4}\label{u3:q7}
Para os inteiros $x$ com $-3\le x\le3$:
\begin{enumerate}[label=\alph*)]
  \item para quais é verdadeira a proposição \enquote{se $x^2>4$, então $x>2$}?
  \item escreva a contrapositiva dessa proposição e verifique que ela é verdadeira para os mesmos valores;
  \item a proposição do item (a) é verdadeira para \emph{todo} inteiro? Justifique.
\end{enumerate}
\resp{a) $x=-2,-1,0,1,2,3$\ \ b) \enquote{se $x\le2$, então $x^2\le4$}; mesmos valores\ \ c) Não: falha em $x=-3$.}
\end{questao}

\begin{questao}{2}{autoral, no formato de concursos}{3-U3, 3-E1}\label{u3:q8}
Sabe-se que é \textbf{falsa} a afirmação: \enquote{Se Ana viajou, então Bruno ficou em casa ou Carla saiu}. Pode-se concluir que:
\begin{enumerate}[label=(\alph*)]
  \item Ana não viajou.
  \item Ana viajou, Bruno ficou em casa e Carla não saiu.
  \item Ana viajou, Bruno não ficou em casa e Carla não saiu.
  \item Ana não viajou ou Bruno ficou em casa.
  \item Nada se pode concluir sobre Bruno.
\end{enumerate}
\resp{(c)}
\end{questao}

\begin{questao}{2}{autoral — demonstração}{3-H3, 3-U4}\label{u3:q9}
\begin{enumerate}[label=\alph*)]
  \item Usando a tabela-verdade, demonstre que $p\rightarrow q$ e ${\nao}q\rightarrow{\nao}p$ têm sempre o mesmo valor, e que $p\rightarrow q$ e $q\rightarrow p$ nem sempre têm.
  \item Use o item (a) para demonstrar que, para todo inteiro $n$, se $n^2$ é múltiplo de 3, então $n$ é múltiplo de 3. (Sugestão: um inteiro que não é múltiplo de 3 é da forma $3k+1$ ou $3k+2$.)
\end{enumerate}
\resp{Demonstração. A prova deve mostrar: a) que as colunas de $p\rightarrow q$ e ${\nao}q\rightarrow{\nao}p$ coincidem nas quatro linhas e que a de $q\rightarrow p$ difere em alguma; b) que, se $n=3k+1$ ou $n=3k+2$, então $n^2$ deixa resto 1 na divisão por 3 ($n^2=3(3k^2+2k)+1$ ou $3(3k^2+4k+1)+1$), logo não é múltiplo de 3 — e que isso, pela contrapositiva, prova o enunciado.}
\end{questao}

\begin{questao}{2}{autoral}{3-U4, 3-E2}\label{u3:q10}
Quatro cartões têm um número de um lado e uma cor do outro. Você vê: \mbox{\fbox{3}\ \fbox{8}\ \fbox{vermelho}\ \fbox{azul}}. Regra: \enquote{se um cartão tem número par, então o outro lado é vermelho}. Quais cartões é preciso virar, e só esses, para verificar a regra?
\resp{O 8 e o azul.}
\end{questao}

\begin{questao}{2}{Simulado OBMEP 2017, nível 3, q.~14 (via FME)}{3-U2}\label{u3:q11}
Ari, Bruna e Carlos almoçam juntos todos os dias e cada um deles pede água ou suco.
\begin{itemize}
  \item Se Ari pede a mesma bebida que Carlos, então Bruna pede água.
  \item Se Ari pede uma bebida diferente da de Bruna, então Carlos pede suco.
  \item Se Bruna pede uma bebida diferente da de Carlos, então Ari pede água.
  \item Apenas um deles sempre pede a mesma bebida.
\end{itemize}
Quem pede sempre a mesma bebida e que bebida é essa?
\begin{multicols}{3}
\begin{enumerate}[label=(\alph*)]
  \item Ari; água
  \item Bruna; água
  \item Carlos; suco
  \item Ari; suco
  \item Bruna; suco
\end{enumerate}
\end{multicols}
\resp{(a)}
\end{questao}

\begin{questao}{3}{autoral — demonstração}{3-U4, 3-H3}\label{u3:q12}
Para saber se 97 é primo, ninguém testa todos os divisores até 96: testam-se só os primos até $\sqrt{97}$, isto é, 2, 3, 5 e 7. Por que isso basta? A justificativa é um condicional: \enquote{se um natural $n>1$ não tem divisor primo $p$ com $p\le\sqrt n$, então $n$ é primo}.
\begin{enumerate}[label=\alph*)]
  \item Escreva a contrapositiva desse condicional.
  \item Demonstre a contrapositiva. (Sugestão: se $n$ não é primo, $n=ab$ com $1<a\le b$.)
  \item Use o teste para decidir se 97 e 91 são primos.
\end{enumerate}
\resp{a) \enquote{se $n>1$ não é primo, então $n$ tem um divisor primo $p\le\sqrt n$}\ \ b) Demonstração. A prova deve mostrar que, com $n=ab$ e $1<a\le b$, vale $a^2\le ab=n$, logo $a\le\sqrt n$, e que um divisor primo de $a$ é $\le a\le\sqrt n$ e divide $n$\ \ c) 97 é primo (não é divisível por 2, 3, 5 nem 7); 91 não é ($91=7\cdot13$).}
\end{questao}

\begin{questao}{3}{adaptada do Simulado OBMEP 2017, nível 2, q.~12}{3-U3, 3-E1}\label{u3:q13}
Tia Geralda sabe que um, e só um, de seus sobrinhos — Ana, Bruno, Cecília, Daniela e Eduardo — comeu todos os biscoitos. Ela sabe também que o culpado sempre mente e que os inocentes sempre dizem a verdade. Três deles falam:
\begin{itemize}
  \item Bruno: \enquote{O culpado é Eduardo ou Daniela.}
  \item Eduardo: \enquote{O culpado é uma menina.}
  \item Daniela: \enquote{Se Bruno é culpado, então Cecília é inocente.}
\end{itemize}
Quem comeu os biscoitos? Mostre que a resposta é única.
\resp{Eduardo.}
\end{questao}

\begin{sintese}
\textbf{Ao fim desta unidade, você deveria saber:}
\begin{itemize}[leftmargin=1.4em,itemsep=1pt]
  \item dar o valor de $p\rightarrow q$ e de $p\leftrightarrow q$ em todos os casos, inclusive com antecedente falso;
  \item traduzir \enquote{suficiente}, \enquote{necessário} e \enquote{somente se} para um condicional;
  \item distinguir contrapositiva (equivalente) de recíproca e inversa (não equivalentes), e usar a contrapositiva numa prova.
\end{itemize}
\smallskip\noindent\textbf{Responda com suas palavras, sem consultar.} Depois, confira com a caixa \emph{Para lembrar} (p.~\pageref{lembrar-u3}) e com a seção indicada.
\begin{enumerate}[label=\arabic*.,leftmargin=1.4em]
  \perguntasintese{O que significa \enquote{$p\rightarrow q$ é verdadeiro}? Em que único caso ele é falso? \hfill\mbox{\small\color{cinza}(→ Definição)}}
  \perguntasintese{Como a promessa explica a tabela de $\rightarrow$, e por que o caso \enquote{antecedente F} tem de ser V? \hfill\mbox{\small\color{cinza}(→ De onde vem?)}}
  \perguntasintese{Como você reconhece, num enunciado, um condicional escondido (\enquote{basta}, \enquote{só se}, \enquote{é preciso})? \hfill\mbox{\small\color{cinza}(→ Onde é usado?)}}
  \perguntasintese{Qual é o erro mais perigoso ao trabalhar com um \enquote{se\ldots\ então}? Como você o detecta? \hfill\mbox{\small\color{cinza}(→ Erros comuns)}}
\end{enumerate}
\end{sintese}

% ------------------------------------------------------------
\Respostas
{\small
\resposta{S}{Síntese de domínio: confira com a caixa \emph{Para lembrar} (p.~\pageref{lembrar-u3}) e com a seção indicada em cada pergunta.}
}
